% SGA 3, Expose V, Loop 1 English reconstruction.
% Scope: completion of Lemma 6.1, proofs (c) and (d).
% Authority: Polo--Gille Exp5-13oct24.pdf, component-local p. 23
% (combined-reader p. 285), re-edited page header 273.
% Jacob Reinhold's CC BY 4.0 Markdown is comparison/drafting material only.

设 \(Y'\) 是 \(Y\) 的开集, 并置
\[
  U'=q^{-1}(Y'),\qquad X'_0=p^{-1}(Y'),
\]
等等.\footnote{\label{note:sga3-V-33}编者注: 下面的论证已扩充,
特别是对 \(U'\) 为 \(X'_0\) 上所诱导群胚的拟截面这一点的验证.}
则 \(U'\) 是 \(U\) 的开集, 对群胚 \(U_\ast\) 所定义的等价关系饱和.
由\hyperref[lem:sga3-V-1.1]{引理~1.1} 和
\hyperref[sga3:exp5:lemma:1.2]{1.2} 可知, \(Y'\)
既在 \((\mathbf{Sch}/S)\) 中, 也在赋环空间范畴中,
是由 \(U_\ast\) 在 \(U'\) 上诱导的群胚的余核.
同样, \(X'_0\) 是 \(X_0\) 的开集, 对 \(X_\ast\) 所定义的等价关系饱和,
并有如下交换图, 其中两个方块都是笛卡儿图:
\SGAThreeDiagram[0.58\textwidth]
  {figures/exp5/p273_lemma61_restricted_groupoids_cartesian_corrected.png}
  {\(X'_0\) 与 \(X_0\) 上的限制群胚.}
  {fig:sga3-V-6.1-restricted-groupoids}
  {第 V 篇, 分篇第~23 页 / 合订本第~285 页 / 重版本页眉页码~273;
  右下角的 \(U\) 由源 PDF 的 \(U'\) 改正}
这里 \(\widetilde d_1\) 满射且有限局部自由.
另一方面, 设 \(x\in U'\). 由于 \(U\) 是拟截面, 集合
\[
  E=d_0d_1^{-1}(x)\cap U
\]
有限, 并包含在 \(U\) 的一个仿射开集 \(W\) 中.
于是 \(E'=E\cap U'\) 是包含在拟仿射开集 \(W\cap U'\) 中的有限集.
因此 \(W\cap U'\) 有一个包含 \(E'\) 的仿射开集 \(W'\).
这就证明了 \(U'\) 是由 \(X_\ast\) 在 \(X'_0\) 上诱导的
群胚 \(X'_\ast\) 的拟截面.
将\hyperref[lem:sga3-V-6.1-i]{(i) 项}的第一个断言应用于
\(X'_\ast\) 和 \(U'\), 即知 \(Y'\) 是 \((\mathbf{Sch}/S)\)
中 \(X'_\ast\) 的余核.

特别地, 对每个 \(S\)-概形 \(T\), 序列
\begin{equation}
  T(Y')
  \xrightarrow{\,T(p|_{X'_0})\,}
  T(X'_0)
  \underset{T(d_0|_{X'_1})}
    {\overset{T(d_1|_{X'_1})}{\rightrightarrows}}
  T(X'_1)
  \label{eq:sga3-V-6.1-T-exact}
\end{equation}
正合. 若 \(T\) 是仿射直线 \(\mathbf G_{a,S}\)
(\SGARefLink{sga3:i:subsection:4-3}{第 I 篇, 4.3}), 则此序列与
\begin{equation}
\begin{aligned}
  \Gamma(Y',\mathcal O_Y)
  &\longrightarrow
  \Gamma(p^{-1}(Y'),\mathcal O_{X_0})
  \underset{\delta_0}{\overset{\delta_1}{\rightrightarrows}}
  \Gamma(d_0^{-1}p^{-1}(Y'),\mathcal O_{X_1})\\
  &\hspace{9em}=
  \Gamma(d_1^{-1}p^{-1}(Y'),\mathcal O_{X_1}),
\end{aligned}
\label{eq:sga3-V-6.1-sections-exact}
\end{equation}
等同, 因而后者对每个开集 \(Y'\) 都正合.
这就完成了引理~\hyperref[lem:sga3-V-6.1-i]{6.1(i)} 的证明.

\medskip
\phantomsection\label{par:sga3-V-6-d}%
\noindent\textbf{d) \hyperref[lem:sga3-V-6.1-iv]{(iv)} 的证明.}---%
若 \((d_0,d_1)\) 是等价对, 则 \((u_0,u_1)\) 也是.
由定理~\hyperref[thm:sga3-V-4.1]{4.1} 可知,
\[
  u_0\boxtimes u_1:U_1\longrightarrow U\times_Y U
\]
是同构, 因而 \(v_0\boxtimes v_1\) 也是同构;
见步骤~\hyperref[par:sga3-V-6-b]{(b)} 中的笛卡儿方块.
由于 \(v\times v\) 忠实平坦且拟紧,
\(d_0\boxtimes d_1\) 是同构 (SGA~1, VIII, \SGARefLink{sga3:viii:corollary:5:4}{5.4}).

此外, 若 \(d_0\) 平坦, 则 \(u\) 平坦.
由定理~\hyperref[thm:sga3-V-4.1]{4.1}, 态射 \(q\) 平坦, 故 \(r\) 也平坦.
由于 \(v\) 忠实平坦, \(p\) 平坦; 它还满射, 因而忠实平坦.

% Translation coverage: entire assigned component; adjoining components retain the original statement/proof seams.
% Source lines: 86; translated body lines before this audit comment: 82.
% Math preserved: 53 inline formulas, 3 display blocks, 2 equation environments; mathematical prose translated where present.
% Labels: 4/4; hyperrefs: 8/8; SGARefLinks: 2/2; inline targets: 0/0; footnotes: 1/1.
% Figure arguments preserved; all rendered captions and natural-language source locators translated.
% Unresolved TODOs: 0; compilation and visual inspection remain with the volume integrator.
